% !TeX root = ../papers/nc-p3.tex
\section{外側の極における二つの小さい変数}
\label{sec:rank-fifteen-outer-pole}
\begin{lemma}\label{lem:rank-fifteen-outer-pole}
剰余標数3で $0<2\nu(x-1)<\nu(y)$ とする。付値を延長すると、
$\nu(w)=\nu(x-1)/2$, $\nu(u)=(\nu(y)-2\nu(x-1))/4$ の表示が存在する。
$H=uw$, $D=(1-2H^2)^{-1}$, $T=1-2w^2$ とすると
\[
             x=(1+2H^2D^2)T^2,\qquad y=H^4D^4 T^6.
\]
経路係数は $(u,w)=(0,0)$ における $\mathbb Z_{(3)}$ 上の整な形式級数で表せる。
\end{lemma}
\begin{proof}
補題\ref{lem:rank-fifteen-ramified-pole}と同じ $r,K,h,H$ を選ぶと
$\nu(H)=\nu(y)/4$ である。同補題の $1-2AB$ の恒等式では
$\nu(r^2)>\nu(x-1)$ なので $w^2=(1-2AB)/2$ は
$\nu(w^2)=\nu(x-1)$ を満たす。$u=H/w$ は正付値である。
そこで同補題の $HL=w$, $HC=w\sqrt{1-u^2}$ を選び、
$E=\sqrt{1-w^2}$, $F=\sqrt{1-w^2-H^2}$ を代入する。
各平方根の定数項は1または2であり、2は単元なので全係数は上記の環にある。
この代入で $x,y$ の二式も得る。
\end{proof}

\begin{theorem}\label{thm:rank-fifteen-outer-pole}
標数零の通常な階数15窓で、剰余標数3において $0<4\nu(x-1)=\nu(y)$ とする。
この窓は標準Koszul左右AS正則で、各始点で四中心四次式上自由階数256である。
\end{theorem}
\begin{proof}
前補題で $\nu(u)=\nu(w)$ となるので、$s=w^2$, $m=u/w$ とすれば $H=ms$ で
$m$ は単元である。$X=(1+2m^2s^2D^2)T^2$, $Y=m^4s^4D^4T^6$ と置く。
$B=(X-1)/s$ は整で $B(0,m)=-4$ であり、
$\Lambda=Y/(X-1)^4=m^4D^4T^6/B^4$ は $s=0$ で $m^4/256$ に等しい。
$\partial_s X(0,m)=-4$, $\partial_m X(0,m)=0$,
$\partial_m\Lambda(0,m)=m^3/64$ より、どの非零な剰余根でもヤコビ行列は可逆である。
従って $\bar m^4=256\bar\Lambda$ の四根のうち任意のものを指定して
所定の $x,y$ に持ち上げることができる。

\path{computations/verify_rank_fifteen_outer_pole.py} は $\mathbb Q(m)$ 上の
三十二桁を用いる。$m-1$ と3の上で単元の主元を選ぶと、未切断経路での最大損失は
$(0,3,3,4,4)$、四次中心の各文字では $(0,2,2,2)$、四次・五次の積表では
$(0,3)$ である。長さ6から13の $(3,1,2,1,3,0,2,0)$ 桁を使って最後の一桁で
十三次の消滅を得る。各 $w$ 除算前の定数項は有理関数として零であり、
使用した定数項の分子・分母の剰余の平方自由な積は
\[
                    (m^9-m)/(m-1)
\]
である。従って $\bar m\notin\mathbb F_9$ または $\bar m=1$ ならこの表示が使える。
$\mathbb F_9^\times$ の四乗は $\{1,-1\}$ である。四乗が1なら $\bar m=1$ を選べる。
四乗が $-1$ なら $m^2+m+2$ の根を選ぶ。この既約式と3の上で単元の主元を取り直し、
同じ三十二桁の標数零計算を十三次まで行うと、用いた全ての主元の分子・分母は
剰余でこの既約式と互いに素である。全ての剰余値で中心商の上界と一般の下界が一致し、
第\ref{sec:rank-fifteen-overlap-annulus}節の判定が適用できる。
\end{proof}
